\section*{Problem Set 5 due on April 11 at 11:59pm}

\question{1}{}
In the Figure below 

\begin{itemize}
    \item [(a)]In the upper left corner of the figure, there are no solutions where the Universe is currently expanding from a singular state. Derive the analytic form for the boundary of the shaded region.
    \item [(b)] The curve above the line $\Omega_\Lambda = 0$ separates universes which either expand for ever or recollapse. Derive the analytic form for this curve.
    \item [(c)] In the lower right, the shaded region corresponds to an exclusion due to the age of the Universe. Assuming $H_0 = 71$ km/s/Mpc calculate the age of the Universe along the boundary of the region.
\end{itemize}


\answer{}



\clearpage
\question{2}{}
Find the proper distance and determine the time the light was emitted
\begin{itemize}
    \item [(a)] to a galaxy with redshift $z=0.01$
    \item [(b)] to a galaxy with redshift $z=3$
    \item [(c)] What is the angle spanned by the horizon at a redshift $z=1100$? You will need to know the size of the horizon at that redshift and the (angular) distance to that redshift.
\end{itemize}
We will assume matter $\gamma=1$ and compare the results for (i) $(\Omega_0,\Omega_\Lambda)=(1,0)$ and (ii) $(\Omega_0,\Omega_\Lambda)=(0.27,0.73)$.

\answer{}




\clearpage
\question{3}{}
A particle $\chi$ with mass $m=1$ GeV and $g_\chi=2$ as a cross-section for the process $\chi\bar{\chi}\to e^+e^-$ given by
\begin{align}
    \langle \sigma v\rangle =\frac{T^2}{(T^2+m^2)^2}.
\end{align}
Between what temperatures is $\chi$ in thermal equilibrium with the electron? 
\answer{}